\section{Introduction}

Large vision--language models (LVLMs) such as LLaVA-1.5 \cite{liu2023llava} produce fluent image descriptions, but they often \emph{hallucinate}: they mention objects that are not present in the image. Because hallucination is a grounding failure, a natural response has been to inspect or manipulate the model's visual attention. Detection methods score how much, or where, an object token attends over the image \cite{jiang2024devils,hoangxuan2026pas}; mitigation methods amplify image attention or redistribute attention away from sinks \cite{liu2024pai,yin2025clearsight,kang2025visualsink}. These approaches share an appealing premise: better visual attention should imply better visual grounding.

This paper studies a gap in that premise through a verification-centered diagnostic framework: \textbf{attention allocation is not evidence verification}. Attention tells us where computation is routed, but object hallucination depends on whether a specific object claim is supported by visual evidence. The distinction matters in two directions. A detector can assign high hallucination scores because attention statistics correlate with a nuisance variable rather than grounding. A mitigation method can make the model look more at the image, yet still fail to distinguish whether the queried object is present or absent.

A central mechanism is what we call \emph{associated-evidence grounding}. An LVLM may attend to a real visual region that is related to the queried object, while the queried object itself is absent. For example, when asked whether an image contains a \emph{sink}, a model may route attention to a visible \emph{toilet} or bathroom context, activate co-occurrence and embedding-space associations, and answer ``yes.'' The attention map is not random: it points to meaningful visual evidence. But the evidence is not \emph{target-discriminative}. This mechanism is consistent with analyses that connect object hallucination to overconfidence in irrelevant visual features after visual tokens are projected into the language embedding space \cite{bai2025adavib}. It clarifies why attention-based methods can fail even when the model is genuinely looking at the image: a visually coherent attention map can certify associated evidence, but only target-discriminative evidence can certify the object claim.

We first expose this problem on the detection side through \textbf{generation-position confounding}. Captions are produced autoregressively, and the position of an object token affects both the feature and the label. As decoding proceeds, attention shifts from image tokens toward generated text; at the same time, hallucinations become more frequent late in the caption. On our LLaVA-1.5-7B / MSCOCO run, the hallucination rate rises from $2.2\%$ in the first ten generated tokens to over $60\%$ past position 100. A score that tracks position can therefore appear to detect hallucination without testing visual support.

This confound is strong enough to reorder the field's apparent progress. A trivial detector using only generation position reaches $0.830$ AUROC, comparable to or higher than attention baselines. Under our position-controlled protocol---within-bin, matched-pair, and residual AUROC---PAS and SVAR drop from $\sim$$0.83$ overall AUROC to $\sim$$0.58$ within-bin AUROC, retaining only $22$--$28\%$ of their above-chance signal. Representation and uncertainty baselines such as Internal Confidence (IC), entropy, and NLL retain substantially more signal. Even SinkDetect, our stress-test detector built to use purified visual-attention shape rather than mass, collapses to chance under position control.

We then ask whether the same proxy failure appears in \emph{attention-based mitigation}. If visual attention were a reliable grounding handle, amplifying or re-routing it should improve hallucination behavior, not merely change output tendencies. We reproduce the attention-intervention components of PAI, ClearSight, and Visual Attention Sink on POPE and CHAIR. The results are cautionary: PAI-attention-only gives essentially no benefit; ClearSight increases POPE recall but also increases false positives by a comparable or larger amount; and Visual Attention Sink makes captions longer, mentions more objects, and hallucinates more object mentions. These failures are consistent with the same mechanism: increasing visual influence can amplify associated evidence, visual plausibility, or object richness without improving object-presence verification.

The same diagnostic also gives a constructive path forward. We formalize this direction as Target-Discriminative Evidence Verification (TDEV): a verifier should compare target evidence against associated evidence rather than only measuring visual routing. When we instantiate TDEV with independent target-region evidence from OWLv2 \cite{minderer2023scaling}, target absence remains strongly predictive under the same position and same-object controls that defeat attention mass. On POPE, a hybrid gate-plus-rescue rule lowers related-present false positives from $0.114$ to $0.069$ and improves macro MCC from $0.730$ to $0.763$, while still leaving a modest, incomplete mitigation effect. This positive control narrows the conclusion: visual evidence is useful, but the useful signal is target verification rather than undifferentiated visual routing.

Studying detection and mitigation together is important because the two literatures often validate the same proxy in opposite directions. A detector treats an attention pattern as a measurement of grounding; an intervention treats a modified attention pattern as a mechanism for improving grounding. If the proxy is valid, the same signal should remain selective under nuisance controls and should improve object-presence discrimination when amplified. If it is invalid, both tasks can show attractive aggregate gains for the wrong reason. This paper uses that symmetry as a diagnostic tool: detection experiments reveal whether attention scores measure object evidence, mitigation experiments reveal whether attention changes cause object verification, and associated-evidence audits test whether plausible visual evidence is actually target-discriminative.

\paragraph{Contributions.} We make seven contributions.
\begin{itemize}
\item We identify generation-position confounding in token-level LVLM hallucination detection and show that position alone reaches $0.830$ AUROC on LLaVA-1.5-7B / MSCOCO (\S\ref{sec:protocol}).
\item We introduce a verification-centered diagnostic framework and instantiate it with position-controlled detection metrics, answer-prior mitigation metrics, attention-shift selectivity audits, and semantic-neighbor negative slices (\S\ref{sec:protocol},~\S\ref{sec:associated-audit}).
\item We re-evaluate detection baselines spanning uncertainty, representation, and attention signals (\S\ref{sec:protocol},~\S\ref{sec:results}).
\item We build SinkDetect as a stress test of carefully engineered attention-shape features and show that even purified attention shape does not yield position-independent detection signal in this setting (\S\ref{sec:sinkdetect}).
\item We extend the critique to attention-based mitigation by evaluating PAI-attention-only, ClearSight, and Visual Attention Sink under POPE yes-prior diagnostics and CHAIR caption statistics (\S\ref{sec:mitigation}). The results support our central claim: looking is not verifying.
\item We formalize Target-Discriminative Evidence Verification (TDEV) and evaluate an OWLv2 instantiation as a positive control, showing that target-discriminative region evidence survives the detection controls and modestly reduces semantic-neighbor false positives under a hybrid POPE gate-plus-rescue rule (\S\ref{sec:tdev},~\S\ref{sec:results},~\S\ref{sec:associated-audit}).
\item We identify associated-evidence grounding as a mechanism that links the detection and mitigation failures: attention can select visually meaningful but non-discriminative evidence, so future methods must test target-object verification rather than visual routing alone.
\end{itemize}
